% 前言

\section*{前言}
\addcontentsline{toc}{section}{前言}

这些年在课堂上，学生常问我：“云计算究竟要学到什么程度，才算真正入门？”产生这种困惑并不奇怪。
云计算涉及的技术层次多，Docker、Kubernetes、Raft、Serverless、eBPF 等概念又不断出现。面对越来越长的学习清单，学生往往难以判断哪些内容必须掌握，也不知道自己是否真正理解了云计算。

我认为学习云计算的关键，不是继续积累技术名词，而是建立一条由真实问题驱动的学习线索。
沿着这条线索，就可以理解一项技术为什么出现、受到哪些约束、解决了什么问题，又付出了什么代价。
本书希望帮助读者建立的，正是这样一条从问题走向机制、再回到工程取舍的线索。

\subsection*{为什么写这本书}

要建立这条线索，只掌握概念或只熟悉操作都不够。读者常见的云计算学习材料大致从两个方向展开。一类侧重概念体系，帮助读者理解 IaaS、PaaS、SaaS、CAP 等基本概念；另一类侧重平台操作，通过控制台、配置文件和部署步骤帮助读者快速上手。前者便于建立知识框架，后者便于熟悉具体工具，但从理解概念到独立分析和构建一个在线系统，中间仍有一段路需要补足。

云计算的难点通常不是某一个知识点格外陌生，而是许多已经学过的知识会同时作用于一个系统。学过 TCP 三次握手，不等于知道大量长连接到来后为什么会排队；学过进程、线程和锁，不等于知道并发为什么既能提高吞吐，也可能破坏共享状态；了解一致性与可用性的概念，也不等于能够判断一次跨节点操作应当怎样确认结果。网络、操作系统、数据结构和分布式系统原本分散在不同课程中，进入真实在线系统后却必须共同工作。

本书位于这些基础课程与复杂系统实践之间。它从系统中已经出现的问题出发，分析问题受到哪些约束，引入相应的机制，再通过实现和运行结果检验设计是否成立。读者需要掌握的不只是某种工具的用法，还包括机制的适用条件，以及性能、状态正确性、可扩展性、可用性和成本之间的取舍。

要把这些内容讲清楚，需要一个能够观察、实现和验证的教学载体。这些年，我先后尝试过选课系统和大模型服务系统。选课系统贴近生活，但问题主要集中在数据访问和名额分配，难以自然覆盖跨节点协作、故障恢复和弹性伸缩；大模型服务涉及队列调度、资源管理和横向扩展，但依赖 GPU，推理过程也具有较强的黑盒属性，不适合所有学生完整实现。

经过比较，我们最终选择了在线游戏。游戏服务既包含长连接、并发处理、共享状态和跨节点协作，也会遇到部署、故障和弹性伸缩问题；更重要的是，它的运行过程足够透明。一次操作从客户端发出，经过网络传输、服务器校验、状态更新和结果同步，每一步都可以通过代码实现，也可以通过实验观察。学生可以运行配套系统，并直接感受延迟、状态冲突和故障对服务的影响。

从选课系统到大模型服务，再到在线游戏，这条教学路线经过了近十年的调整和打磨。本书是这段课堂实践的阶段性总结。

\subsection*{这本书有什么不一样}

\textbf{第一，以持续演进的在线系统引出技术。}

全书先建立在线系统的整体视图，再从最小的双人在线交互出发，推动案例随着规模和压力逐步演进：建立网络闭环后，处理单机并发；单机容量接近上限后，将系统拆分到多台机器；服务和实例难以人工管理后，再引入容器、编排和弹性机制；最后进入基础设施内部，解释容器隔离、网络虚拟化和弹性调度如何实现。

每一阶段都承接前一阶段已经暴露的问题。读者由此看到的不是一组并列的技术名词，而是问题、约束、机制、取舍和新问题之间的因果关系。

\textbf{第二，把分散的知识组织成完整的系统判断。}

设计消息协议时，TCP 和 UDP 不再只是网络课程中的定义；组织并发时，进程、线程、goroutine、锁和 I/O 多路复用需要落实到真实代码；拆分系统时，哈希、分片、复制、事务和共识会共同决定请求与状态应当由谁处理和保存。

在线游戏是本书的主要教学案例，但不限定知识范围。正文还会引入在线协作、电商平台和大模型服务等系统，说明传输、并发、一致性、容器和调度机制在不同场景中的作用与适用条件。业务形式可以变化，规模、状态、故障和资源问题具有共通性。

\textbf{第三，使原理、实现和验证相互对应。}

正文解释机制及其适用条件，配套案例提供可运行的代码，四组实验则围绕同一个游戏系统逐步扩大规模。读者需要观察吞吐、延迟、资源占用和错误现象，修改其中一个机制，再用运行结果判断修改是否有效。

最后一章还将这些方法延伸到需要大量隔离执行环境的 Agent Sandbox。当平台管理的不再只是服务副本，还包括需要暂停、恢复、回滚和复现的执行过程时，云计算中的隔离、网络、调度、状态和验证机制会以新的方式组合。这个问题既是对全书知识的综合应用，也为希望继续深入的读者提供了进入系统研究的起点。

本书可以用于云计算和系统工程类课程，也可以作为课程设计、综合实践以及进一步学习云原生系统的基础。

\subsection*{这本书如何学}

阅读本书不要求预先掌握云计算或分布式系统，但需要具备基本的编程能力，并了解常见的数据结构。如果已经学过计算机网络和操作系统，阅读会更加顺畅；即使部分知识已经遗忘，也不必先把所有先修课程重新学习一遍，可以跟随案例，在真正需要时重新理解它们。

全书六章构成一条连续的学习路线：

\begin{itemize}
    \item \textbf{第一章“谋定全局：在线系统架构”}建立在线系统的整体视图，并从性能、状态正确性、可扩展性、可用性和成本五个维度评价架构；
    \item \textbf{第二章“双雄对战：网络通信”}建立最小网络闭环，讨论通信、消息组织、状态同步和结果确认；
    \item \textbf{第三章“英雄集结：单机并发”}进入单机并发，处理执行组织、共享状态保护和资源复用；
    \item \textbf{第四章“切分世界：分布式系统”}将系统扩展到多台机器，围绕均衡、协同和容错展开；
    \item \textbf{第五章“飞升入定：云原生部署”}将系统交给集群平台，讨论容器化交付、编排、弹性和自愈；
    \item \textbf{第六章“穷理尽微：云原生核心原理”}继续下沉到容器、网络和弹性调度的内部机制，并探索 Agent Sandbox。
\end{itemize}

学习本书时，有三点需要注意。

\textbf{第一，动手。}配套程序应当真正运行起来。只有观察到连接排队、状态冲突、尾延迟上升或扩容滞后，书中的机制才会从文字变成可以分析的系统行为。

\textbf{第二，追问。}读到一个解决方案时，可以先思考自己会怎样处理当前问题，再把自己的方案与书中的实现进行比较。两者之间的差异，往往能够揭示此前忽略的条件和代价。

\textbf{第三，善用 AI，但不依赖 AI。}AI 可以帮助解释概念、检查代码和讨论替代方案，但不能代替对代码、运行现象和实验结果的判断。能够辨别一个答案是否符合系统的实际约束，本身就是本书希望培养的能力。

对于采用本书授课的教师，本书可以按一学期组织教学。四组实验分别对应网络闭环、单机并发、多节点协作，以及 Kubernetes 上的部署与弹性运行。课时有限时，可以将第六章的拓展内容压缩为导览，但应为前面的系统演进和动手实验保留足够时间。

配套代码、实验材料、授课视频和 PPT 将在 Gitee 仓库持续更新（\url{https://gitee.com/hnu-cloudcomputing}）。实践时应使用相互对应的书稿和代码版本，避免接口或环境差异影响实验结果。

云计算的水面之上，新技术来了又去；水面之下，进程、网络、状态、调度和故障这些基本问题长期存在。理解这些问题，才能在工具变化时保持判断。愿读者既能深入机制，也能回到系统全局。

\subsection*{致谢}

本书由陈果、徐方林负责编写。各章初稿撰写分工如下：胡文举（第1章）、庞海鑫（第2章）、谢先衍（第3章）、贺臻（第4章）、张道平（第5章）、徐方林（第6章）。胡文举承担了课程实验内容的验证与复核，并参与指导了书中配图的视觉修订工作；终稿由徐方林统一核定。

实验案例的设计与实现离不开历届学生的参与。陈俊杰、王煌等同学在游戏案例代码、实验文档和测试用例方面做了大量工作。课程视频的制作同样有赖于学生团队的协助，刘筱芊等同学参与了授课视频的录制、剪辑与后期整理。

出版社的编辑老师们为本书的出版同样付出了大量心血，在此一并感谢。

受知识水平和时间精力所限，本书难免存在错误和疏漏。欢迎读者将发现的问题反馈至邮箱 \texttt{guochen@hnu.edu.cn}，我们将持续改进本书的质量。

\vspace{2em}

\begin{flushright}
陈果\\
2026年7月于长沙
\end{flushright}
